% \documentclass{cumcmthesis}
\documentclass[withoutpreface,bwprint]{cumcmthesis} %去掉封面与编号页，电子版提交的时候使用。


\usepackage[framemethod=TikZ]{mdframed}
\usepackage{url}   % 网页链接
\usepackage{subcaption} % 子标题
\title{\textbf{注意力机制原理、分类及应用}}
\author{2021113140-符世博}
\date{}



\begin{document}
\maketitle

\section{通俗理解Attention}

人在观察事物时会有选择性的关注较为重要的信息，称其为注意力。
通过持续关注这一关键位置以获得更多的信息，而忽略其他的无用信息，
这种视觉注意力机制大大提高了人类处理信息的效率和准确性。
深度学习中的注意力机制和人类视觉的注意力机制类似，
就是在更多信息中把注意力集中放在重要的点上，选出关键信息，
而忽略其他不重要的信息。

\begin{figure}[H]
    \centering
    \includegraphics[width=.9\textwidth]{1.png}
    \caption{人类在看到一副图像时如何分配注意力}
    \label{fig:1}
\end{figure}

图\ref{fig:1}形象化展示了人类在看到一副图像时如何高效分配有限的注意力资源。其中红色区域表明视觉系统更关注目标，很明显对于上图所示的场景，人们会把注意力更多投入到人的脸部，文本的标题以及文章首句等位置。

注意力机制的灵感来源可以归结到人对环境的生理感知，当人类看东西时，一般会将注意力注视着某个地方，而不会关注全部所有信息。比如，视觉系统更倾向于挑选影像中的部分信息进行集中分析，忽略图像中无关的信息。再比如，当看到图\ref{fig:2}时，主要会将目光停留在猫的脸部，以及注意猫的躯干，而后面的草地则会被当成背景忽略，这说明人类在每一处空间位置上的注意力分布不一样。

\begin{figure}[H]
    \centering
    \includegraphics[width=.9\textwidth]{2.png}
    \caption{猫在草地上}
    \label{fig:2}
\end{figure}

通过这种方式，人类在需要重点关注的目标区域，会投入更多的注意力资源，以获取更多的细节信息，而抑制其他区域信息。这样使人类能够利用有限的注意力资源从大量信息中快速获取高价值的信息，极大地提升了大脑处理信息的效率。

\section{注意力机制原理}

注意力机制(Attention Mechanism)是机器学习中的一种数据处理方法，广泛应用在自然语言处理、图像识别以及语音识别等各种不同类型的机器学习任务中。注意力机制对不同信息的关注程度（重要程度）由权值来体现，注意力机制可以视为查询矩阵(Query)、键(key)以及加权平均值构成了多层感知机(Multilayer Perceptron, MLP)。

注意力的思想，类似于寻址。给定Target中的某个元素Query，通过计算Query和各个Key的相似性或相关性，得到每个Key对应Value的权重系数，然后对Value进行加权求和，即得到最终的Attention数值。所以，本质上Attention机制是Source中元素的Value值进行加权求和，而Query和Key用来计算对应Value的权重系数。

\begin{figure}[H]
    \centering
    \includegraphics[width=.9\textwidth]{3.png}
    \caption{注意力机制}
    \label{fig:3}
\end{figure}

在图\ref{fig:3}中，各个参数定义如下

\begin{itemize}
    \item Source：由一系列的<Key, Value>键值对构成
    \item Query：给定的Target元素
    \item Key：Source中元素的Key值
    \item Value：Source中元素的Value值
    \item 权重系数：Query与key的相似性或相关性，权重系数$Similarity(Query,Key_i)$
    \item Attention Value：对Value值进行加权求和；
\end{itemize}

根据注意力机制的原理可知，其计算公式如下：

\begin{equation}
    Attention(Query,Source)=\sum_{i=1}^{L_x}{Similarity(Query,Key_i)*Value_i}
    \label{eq:1}
\end{equation}

其中，$L_x=||Source||$，表示Source的长度。

Attention从大量信息中有选择地筛选出少量重要信息并聚焦到这些重要信息上，忽略大多不重要的信息。聚焦的过程体现在权重系数的计算上，权重越大，越聚焦在对应的Value值上，即权重代表了信息的重要性，而Value是其对应的信息。

\section{注意力机制计算过程}

大多数方法采用的注意力机制计算过程可以细化为如下三个阶段。

三阶段的注意力机制计算流程：

\begin{itemize}
    \item 第一阶段，计算Query和不同Key的相关性，即计算不同Value值的权重系数
    \item 第二阶段，对上一阶段的输出进行归一化处理，将数值的范围映射到0和1之间
    \item 第三阶段，根据权重系数对Value进行加权求和，从而得到最终的注意力数值
\end{itemize}

\begin{figure}[H]
    \centering
    \includegraphics[width=.9\textwidth]{4.png}
    \caption{注意力机制计算过程}
    \label{fig:4}
\end{figure}

\subsection{第一阶段}

计算 Query和不同 Key 的相似性或者相关性，即计算不同 Value 值的权重系数。

计算两者相似性或者相关性常用方法：

\begin{itemize}
    \item 点积：$Similarity(Query,Key_i)=Query\cdot Key_i$
    \item Cosin相似性：$Similarity(Query,Key_i)=\frac{Query\cdot Key_i}{||Query||\cdot ||Key_i||}$
    \item MLP网络：$Similarity(Query,Key_i)=MLP(Query,Key_i)$
\end{itemize}


\subsection{第二阶段}

对第一阶段得到的原始权重进行归一化处理，将数值的范围映射到 0 和 1 之间。

根据产生方法的不同，第一阶段产生的分值的取值范围也不一样，第二阶段引入类似SoftMax的计算方式对第一阶段的得分就行数值转换。一方面，可以进行归一化，将原始计算分值整理成所有元素权重之和为1的概率分布；另一方面，通过SoftMax的内在机制更加突出重要元素的权重。即一般采用式\ref{eq:2}。

\begin{equation}
    a_i=Softmax(Sim_i)=\frac{e^{Sim_i}}{\sum_{j=1}^{L_x}{e^{Sim_j}}}
    \label{eq:2}
\end{equation}


\subsection{第三阶段}

根据权重系数对Value进行加权求和，从而得到最终的注意力数值：

\begin{equation}
    Attention(Query,Source)=\sum_{i=1}^{L_x}{a_i\times Value_i}
    \label{eq:3}
\end{equation}

\section{注意力机制分类}

\subsection{按Attention的可微性分类}

\subsubsection{硬注意力机制}
硬注意力更加关注点，也就是图像中的每个点都可能延伸出注意力。同时，硬注意力是一个随机的预测过程，更强调动态变化。硬注意力是一个不可微的注意力，训练过程往往是通过增强学习(reinforcement learning) 来完成。

硬注意力是0/1问题，某个区域要么被关注，要么不关注，这是一个不可微的注意力。

\subsubsection{软注意力机制}

软注意力更加关注区域或者通道，软注意力是确定性的注意力，学习完成后可以直接通过网络生成，最关键的地方是软注意力是可微的。可微的注意力可以通过神经网络计算梯度，通过前向传播和后向反馈来学习得到注意力的权重。

软注意力是[0,1]间连续分布问题，用0到1的不同分值表示每个区域被关注的程度高低。

在CV领域中，很多相关工作（例如，分类、检测、分割、生成模型、视频处理等）都在使用软注意力，这些工作也衍生了很多不同的软注意力使用方法。这些方法共同的部分，都是利用相关特征学习权重分布，再用学出来的权重施加在特征之上，进一步提取相关知识。

软注意力施加权重的方法：

\begin{itemize}
    \item 加权可以作用在原图上
    \item 加权可以作用在空间维度上，给不同区域加权
    \item 加权可以作用在channel维度上，给不同通道特征加权
    \item 加权可以作用在不同时刻历史特征上，结合循环结构添加权重。例如，机器翻译，或者视频相关的工作
\end{itemize}

\subsection{按Attention的关注域分类}

\subsubsection{空间域}

关注空间域的注意力机制是指在深度学习中，通过模拟人类注意力机制的方式，对输入数据中的不同空间位置进行加权处理，从而提高模型对重要信息的关注程度，减少对无关信息的干扰，以提升模型性能。

在自然语言处理领域，注意力机制最初是应用于序列到序列模型（Sequence-to-Sequence, Seq2Seq）中，用于对输入序列中不同位置的单词进行加权处理，以便模型能够更好地理解和生成序列数据。后来，注意力机制被引入到计算机视觉领域，用于处理图像数据中的不同区域或像素点，实现对图像内容的关注和理解。

关注空间域的注意力机制通常包括以下步骤：
\begin{enumerate}
    \item 提取特征：首先，通过卷积神经网络（CNN）等方法提取输入数据的特征表示
    \item 计算注意力权重：根据输入数据的特征表示，计算不同空间位置的注意力权重，以确定每个位置对模型输出的贡献程度
    \item 加权汇总：根据计算得到的注意力权重，对输入数据的特征表示进行加权汇总，得到最终的表示结果
    \item 应用于任务：将加权汇总后的特征表示应用于具体的任务，如分类、生成等
    
\end{enumerate}

关注空间域的注意力机制在图像分类、图像分割、目标检测等任务中得到广泛应用，能够帮助模型更好地理解图像内容，提高模型性能和泛化能力。通过引入注意力机制，模型可以有选择地关注图像中的重要区域，从而提高对关键信息的提取和利用效率。


\subsubsection{通道域}

关注通道域的注意力机制是指在深度学习中，通过对输入数据中不同通道的特征进行加权处理，以提高模型对重要通道的关注程度，减少对无关通道的干扰，从而提升模型性能。

在深度神经网络中，特征图通常由多个通道组成，每个通道表示了不同的特征信息。关注通道域的注意力机制旨在通过学习每个通道的权重，使模型能够自适应地选择性地关注重要的通道信息，从而提高模型的表征能力和泛化能力。

关注通道域的注意力机制通常包括以下步骤：

\begin{enumerate}
    \item 提取特征：通过卷积神经网络等方法提取输入数据的特征表示，得到包含多个通道的特征图
    \item 计算通道注意力权重：根据特征图中每个通道的特征表示，计算每个通道的注意力权重，以确定每个通道对模型输出的贡献程度
    \item 加权汇总：根据计算得到的通道注意力权重，对特征图中的每个通道进行加权汇总，得到最终的特征表示
    \item 应用于任务：将加权汇总后的特征表示应用于具体的任务，如分类、目标检测等
\end{enumerate}

关注通道域的注意力机制在图像分类、语义分割、目标检测等任务中得到广泛应用。通过引入通道注意力机制，模型可以自适应地学习每个通道的重要性，从而提高对关键特征的提取和利用效率，进而提高模型性能和泛化能力。通过对不同通道的特征进行加权处理，模型可以更好地适应不同数据集和任务的特点，提高模型的灵活性和泛化能力。

\section{注意力机制的应用}

\subsection{在图像描述领域}

图片描述是深度学习的一个典型应用，即输入一张图片，AI 系统根据图片上的内容输出一句描述文字出来。如图\ref{fig:7}，左边是输入原图，下边的句子是 AI 系统自动生成的描述文字，右边是当 AI 系统生成划横线单词的时候，对应图片中聚焦的位置区域。

\begin{figure}[H]
    \centering
    \includegraphics[width=.9\textwidth]{7.png}
    \caption{图片描述上的注意力机制}
    \label{fig:7}
\end{figure}

\subsection{在机器翻译领域}

在机器翻译领域，注意力机制是一种非常重要的技术，它的应用极大地提高了神经机器翻译模型的性能和效果。

在传统的神经机器翻译模型中，整个输入序列的信息会被编码成一个固定长度的向量，这个向量将作为解码器的输入，用于生成目标语言的翻译。然而，当输入序列很长时，这种固定长度的向量可能无法充分捕捉输入序列的所有信息，导致翻译质量下降。注意力机制的引入可以解决这个问题，它允许模型在翻译过程中动态地关注源语言句子中与当前正在翻译的单词最相关的部分，从而提高翻译的准确性和流畅性。

在机器翻译中，注意力机制的核心思想是在解码器每个时间步都计算一个注意力分布，该分布表示了当前要翻译的目标语言单词与源语言句子中各个位置的相关性。这个注意力分布会被用来加权源语言句子的编码器隐藏状态，从而生成上下文向量，作为解码器的输入。通过这种方式，模型可以动态地关注源语言句子中不同位置的信息，有选择地获取最相关的上下文信息，从而提高翻译的准确性。

在机器翻译中，注意力机制有多种具体应用，包括但不限于：

\begin{itemize}
    \item 编码器-解码器模型中的注意力机制：在编码器-解码器模型中，注意力机制被应用于解码器每个时间步，以帮助模型动态地关注源语言句子中的不同部分
    \item 自注意力机制（Self-Attention）：自注意力机制允许模型在处理长距离依赖关系时更加高效，提高了模型的性能和泛化能力
    \item 多头注意力机制（Multi-Head Attention）：多头注意力机制允许模型同时关注不同的表示子空间，提高了模型的表达能力和学习能力
\end{itemize}

注意力机制在机器翻译中的应用使得模型能够更好地捕捉输入序列的信息，提高了翻译的准确性和流畅性。随着深度学习技术的不断发展，注意力机制在机器翻译领域的应用将会越来越广泛，为机器翻译模型的进一步提升提供了重要的技术支持。

\end{document}